\documentclass[conference]{IEEEtran}
\IEEEoverridecommandlockouts
% The preceding line is only needed to identify funding in the first footnote. If that is unneeded, please comment it out.
\usepackage{amsmath,amssymb,amsfonts}
\usepackage{algorithmic}
\usepackage{graphicx}
\usepackage{textcomp}
\usepackage{xcolor}

\usepackage{booktabs} % For professional looking tables
\usepackage{hyperref} % For clickable links and citations
\usepackage[authoryear]{natbib} % For bibliography management (or use biblatex)
\usepackage{enumitem}   % Allows fine-tuning of lists inside tables to save space

\def\BibTeX{{\rm B\kern-.05em{\sc i\kern-.025em b}\kern-.08em
    T\kern-.1667em\lower.7ex\hbox{E}\kern-.125emX}}
\begin{document}

\title{Sentiment Analysis of Code-Mixed Singlish (Sinhala-English) Reviews with Emoji Fusion NLP Approach\\
}

\author{\IEEEauthorblockN{K. M. Sajana Karunanayake}
    \IEEEauthorblockA{\textit{University of Moratuwa} \\
        \textit{Department of Computer Science and Engineering}\\
        Moratuwa, Sri Lanka \\
        sajanak.26@cse.mrt.ac.lk | karunanayakekmsn.26@uom.lk}
}

\maketitle

\begin{abstract}
    Sentiment analysis of code-mixed Sinhala-English text --- commonly termed ``Singlish'' --- has recently begun to attract research
    attention, with annotated datasets and fine-tuned multilingual transformers (e.g., XLM-R) now reported for domains such
    as banking feedback and e-commerce product reviews \citep{aththanayaka2020sentimental, rizvi2025enhancing, galwatta2026salangabhava, dhananjaya2022bertifying}.
    However, this emerging body of work consistently treats emoji as noise to be stripped during preprocessing rather than
    as a signal to be modeled, and the existing Daraz.lk-based dataset targets aspect-based extraction rather than emoji-fused
    overall sentiment. Separately, emoji-aware sentiment fusion has been explored for other code-mixed language pairs, such as
    English-Hindi \citep{singh2024predicting}, but not for Sinhala-English. This research targets that specific,
    verified gap: it proposes to (i) construct an annotated corpus of Singlish reviews that retains emoji, drawing on
    Daraz.lk product reviews and currently unstudied Sri Lankan mobile app store reviews; (ii) fine-tune multilingual
    transformer models on this corpus, building on established Sinhala transformer benchmarking \citep{dhananjaya2022bertifying};
    and (iii) design and evaluate explicit emoji-text fusion strategies, benchmarked against text-only baselines. The expected
    contributions are a new emoji-preserving annotated dataset spanning e-commerce and app-store reviews, and the first controlled
    emoji-fusion benchmark for Sinhala-English code-mixed sentiment analysis.
\end{abstract}

% \section{Introduction}
% % TODO: Motivate the problem -- rise of code-mixed, emoji-rich user-generated content in Sri Lankan e-commerce/app ecosystems; why sentiment analysis matters for businesses and platforms; why Sinhala/Singlish + emoji is an underserved combination.

% \section{Problem Statement}
Recent work has established that Sinhala-English code-mixed sentiment analysis
is feasible: annotated corpora exist \citep{aththanayaka2020sentimental}, XLM-R
has been benchmarked as the strongest available multilingual encoder for
Sinhala \citep{dhananjaya2022bertifying}, and fine-tuned XLM-R models now
report 88--92\% accuracy on Sinhala and code-mixed banking feedback
\citep{rizvi2025enhancing}. A concurrent aspect-based sentiment dataset built
from Daraz.lk product reviews, SalAngaBhava, further closes the e-commerce
review gap \citep{galwatta2026salangabhava}. What remains unsolved is narrower
than originally assumed: (1) every one of these Sinhala-English studies removes
or ignores emoji during preprocessing, so no existing work measures whether
emoji carry sentiment signal beyond what the code-mixed text already encodes,
or how best to fuse the two; (2) emoji-text fusion architectures have been
demonstrated for other code-mixed pairs, e.g. English-Hindi
\citep{singh2024predicting}, but this design has not been adapted or evaluated
for Sinhala-English; and (3) SalAngaBhava's Daraz.lk coverage is aspect-based
rather than emoji-fused overall sentiment, and Sri Lankan mobile app store
reviews --- an emoji-dense source in practice --- have no published annotated
sentiment data at all. (Ride-hailing platforms such as PickMe were considered
as an additional source but do not expose a public review/comment feature, so
they are excluded from the data collection plan.) This research is scoped to
these three specific gaps rather than to code-mixed Sinhala sentiment analysis
in general.

% \section{Research Objectives}
% % TODO: List 3-5 concrete objectives, e.g.:
% % 1. Construct and publicly release an annotated Singlish-emoji review corpus.
% % 2. Benchmark multilingual transformer models (XLM-R, etc.) on this corpus.
% % 3. Design and evaluate emoji-text fusion strategies.
% % 4. Analyze error patterns specific to code-mixing and emoji ambiguity.

% \section{Literature Review}
% % TODO: Cover (a) Sinhala/Singlish NLP and code-mixed sentiment analysis generally, (b) emoji sentiment/embedding research, (c) multilingual transformer fine-tuning (XLM-R, mBERT), (d) existing low-resource review datasets, and identify the specific gap this thesis fills.

% \section{Methodology}
% \subsection{Data Collection}
% % TODO: Scraping approach for Daraz.lk product reviews and Sri Lankan mobile app store reviews (Google Play/App Store); note PickMe has no public review feature and is excluded; ethical/ToS considerations; target corpus size.
% \subsection{Annotation}
% % TODO: Annotation scheme (e.g., positive/negative/neutral, or fine-grained), number of annotators, inter-annotator agreement metric (e.g., Cohen's/Fleiss' kappa).
% \subsection{Model Development}
% % TODO: Baselines (lexicon-based, classical ML); fine-tuning XLM-R or similar multilingual transformer; emoji representation (emoji2vec, learned embeddings) and fusion strategy (early/late fusion, attention-based).
% \subsection{Evaluation}
% % TODO: Metrics (accuracy, macro-F1), ablations (text-only vs. emoji-fused), error analysis plan.

% \section{Expected Contributions}
% % TODO: (1) Novel public dataset; (2) benchmark results for code-mixed Sinhala sentiment analysis; (3) emoji fusion methodology generalizable to other low-resource code-mixed settings.

\clearpage
\nocite{*}
\bibliographystyle{plainnat} % Or choose your preferred style like apa, ieee, etc.
\bibliography{references} % Points to your references.bib file

\end{document}